% Generated. Do not edit.
% Projection scriptural-date-location/evidence, version 1: the Scriptural Date and
% Location page, generated whole from research/chronology.toml, the Scripture
% chronology corpus and the source library by tools/tpt document-types refresh.
% document: liturgy/roman-rite/1962/propers/temporal/51-eleventh-after-pentecost
\section*{Appendix: Evidence for the Scriptural Dates}
\label{sec:date-location-evidence}
{\small For each passage, as the Scriptural Date and Location table names it: what is dated, how the corpus weighs each answer, the source's own words, and the sources cited. A subject the corpus does not date is reported as such and given no date.\par}
{\small
\par\smallskip\noindent\textit{Gradual: Ps. 27:7 (Vg.; Heb. 28), Ps. 27:1 (Vg.; Heb. 28)}\par\nobreak
\begin{itemize}[leftmargin=1.2em,itemsep=1pt,topsep=1pt,parsep=0pt]
\item \textbf{Historical setting}. No date: the chronology corpus has not yet dated this (research pending).
\item \textbf{Composition} (preferred). Psalms: Before c.~165 B.C. (``before the Maccabean period, around 165 B.C.; no individual psalm can be dated securely''; Catholic critical chronology; New American Bible, Revised Edition --- USCCB web state (2026), Psalms introduction).
\end{itemize}
\par\smallskip\noindent\textit{Offertory: Ps. 29:2--3 (Vg.; Heb. 30)}\par\nobreak
\begin{itemize}[leftmargin=1.2em,itemsep=1pt,topsep=1pt,parsep=0pt]
\item \textbf{Superscription setting} (preferred). The building of David's house at Jerusalem: And they built a house for David (The Holy Bible, Douay-Rheims (Challoner revision), Project Gutenberg eBook 1581: 2 Kings 5:7, 5:11; The Catholic Encyclopedia, Volume 4 (New York, 1908)).
\item \textbf{Composition}. As under Gradual: Ps. 27:7 (Vg.; Heb. 28), Ps. 27:1 (Vg.; Heb. 28).
\end{itemize}
\par\smallskip\noindent\textit{Introit: Ps. 67:6--7, 36 (Vg.; Heb. 68), Ps. 67:2 (Vg.; Heb. 68)}\par\nobreak
\begin{itemize}[leftmargin=1.2em,itemsep=1pt,topsep=1pt,parsep=0pt]
\item \textbf{Historical setting}. As under Gradual: Ps. 27:7 (Vg.; Heb. 28), Ps. 27:1 (Vg.; Heb. 28).
\item \textbf{Composition}. As under Gradual: Ps. 27:7 (Vg.; Heb. 28), Ps. 27:1 (Vg.; Heb. 28).
\item \textbf{Prophetic referent} (disputed). The Nativity of Our Lord: B.C.~3 (``the year of Rome 750 which he styles 3 B.C.''; The Catholic Encyclopedia, Volume 3 (New York, 1908)); c.~B.C.~7 (``probably the year 7 B.C.''; The Catholic Encyclopedia, Volume 3 (New York, 1908)). Place: Bethlehem as the Christian referent (The Holy Bible, Douay-Rheims (Challoner revision), Project Gutenberg eBook 1581: Matt. 2:1).
\end{itemize}
\par\smallskip\noindent\textit{Alleluia: Ps. 80:2--3 (Vg.; Heb. 81)}\par\nobreak
\begin{itemize}[leftmargin=1.2em,itemsep=1pt,topsep=1pt,parsep=0pt]
\item \textbf{Historical setting}. As under Gradual: Ps. 27:7 (Vg.; Heb. 28), Ps. 27:1 (Vg.; Heb. 28).
\item \textbf{Composition}. As under Gradual: Ps. 27:7 (Vg.; Heb. 28), Ps. 27:1 (Vg.; Heb. 28).
\end{itemize}
\par\smallskip\noindent\textit{Communion: Prov. 3:9--10}\par\nobreak
\begin{itemize}[leftmargin=1.2em,itemsep=1pt,topsep=1pt,parsep=0pt]
\item \textbf{Final formation} (preferred). The final editorial formation of the Book of Proverbs: Edited as a whole in the early postexilic period (Catholic critical chronology; New American Bible, Revised Edition --- USCCB web state (2026), Proverbs introduction).
\end{itemize}
\par\smallskip\noindent\textit{Gospel: Mark 7:31--37}\par\nobreak
\begin{itemize}[leftmargin=1.2em,itemsep=1pt,topsep=1pt,parsep=0pt]
\item \textbf{Narrated event}. No date: the chronology corpus has not yet dated this (research pending).
\item \textbf{Composition} (disputed). The Gospel of St Mark: A.D.~38--40 (``the tenth or twelfth year after the Ascension (A.D. 38-40)''; The Catholic Encyclopedia, Volume 9 (1910)); A.D.~40 (The Catholic Encyclopedia, Volume 9 (1910)); A.D.~43 (``the third year of Claudius (A.D. 43)''; The Catholic Encyclopedia, Volume 9 (1910)); A.D.~50--67 (``somewhere between A.D. 50 and 67''; The Catholic Encyclopedia, Volume 9 (1910)); A.D.~60--70 (``between 60 and 70''; The Catholic Encyclopedia, Volume 14 (New York, 1912)); After Matthew, and not as late as the destruction of Jerusalem (De auctore, de tempore compositionis et de historica veritate Evangeliorum secundum Marcum et secundum Lucam (1912), responsa-i-ix). Place: Rome traditionally; Egypt reported by Chrysostom (The Catholic Encyclopedia, Volume 9 (1910)).
\end{itemize}
\par\smallskip\noindent\textit{Epistle: 1 Cor. 15:1--10}\par\nobreak
\begin{itemize}[leftmargin=1.2em,itemsep=1pt,topsep=1pt,parsep=0pt]
\item \textbf{Composition} (disputed). The First Epistle to the Corinthians: A.D.~56 (``(1 and 2 Corinthians; Galatians), 56''; The Catholic Encyclopedia, Volume 11 (New York, 1911)); A.D.~57 (``about Easter A.D. 57''; The Catholic Encyclopedia, Volume 4 (New York, 1908)).
\end{itemize}
\par}
